% Figure: Johnson-Nyquist thermal noise. The open-circuit root-mean-square noise
% voltage across a resistor, sqrt(4 k_B T R B), plotted against resistance for a
% fixed bandwidth at two temperatures. The noise is the electrical signature of the
% equilibrium fluctuations the equipartition theorem predicts, and it sets the
% floor beneath every low-level measurement. Curves are the square-root law; the
% higher temperature lifts the whole floor.
\begin{figure}[htbp]
\centering
\begin{tikzpicture}
\begin{axis}[
  width=11cm, height=7cm,
  xlabel={resistance $R$ (k$\Omega$)},
  ylabel={rms noise voltage $v_n$ ($\mu$V)},
  xmin=0, xmax=1000, ymin=0, ymax=5,
  axis lines=left,
  tick label style={font=\scriptsize},
  label style={font=\small},
  ymajorgrids=true,
  grid style={enggray!22},
  legend style={font=\scriptsize, at={(0.03,0.97)}, anchor=north west, draw=enggray!40},
]
% v_n = sqrt(4 k_B T R B), B = 1e6 Hz (1 MHz). With R in k-ohm (=1e3*x ohm) and
% output in microvolt (=1e6 * v). k_B = 1.381e-23.
% v_n(uV) = 1e6 * sqrt(4 * 1.381e-23 * T * 1e3*x * 1e6)
%         = 1e6 * sqrt(5.524e-14 * T * x) . For T=300: 1e6*sqrt(1.657e-11 * x).
\addplot[engblue, very thick, domain=0:1000, samples=140]
  {1e6*sqrt(5.524e-14*300*x)};
\addlegendentry{$T = 300$ K}
\addplot[engred, very thick, dashed, domain=0:1000, samples=140]
  {1e6*sqrt(5.524e-14*77*x)};
\addlegendentry{$T = 77$ K (liquid nitrogen)}
% mark a 1 M-ohm resistor at room temperature
\addplot[mark=*, only marks, englime!60!black, mark size=2.4pt]
  coordinates {(1000,{1e6*sqrt(5.524e-14*300*1000)})};
\node[font=\scriptsize, anchor=east, englime!45!black] at (axis cs:960,4.3)
  {$1\,\mathrm{M}\Omega$ at $300$ K};
\end{axis}
\end{tikzpicture}
\caption[Johnson-Nyquist thermal noise across a resistor]{The open-circuit
thermal-noise voltage $v_n=\sqrt{4k_B T R B}$ across a resistor, plotted against
resistance for a fixed $1\,\text{MHz}$ bandwidth at room temperature and at liquid-
nitrogen temperature. The noise is the voltage a resistor generates from the random
thermal motion of its charge carriers, the electrical face of the equilibrium
fluctuations equipartition predicts, and its mean-square value $4k_B T R B$ carries
the temperature directly: a resistor is a thermometer that reads itself out as
noise. The square-root law makes the floor grow slowly with resistance, so a
megohm source at room temperature already contributes microvolts across a wide
band. Cooling to $77\,\si{\kelvin}$ drops the floor by the factor
$\sqrt{300/77}\approx2.0$, which is why low-noise front ends for faint signals are
cryogenically cooled. The noise power delivered to a matched load, $k_B T B$, is
independent of resistance, the same $k_B T$ per mode that sets every thermal energy
scale in the chapter.}
\label{fig:vol04ch04:thermal-noise}
\end{figure}
